
\section{Related Works}
\label{sec:related_works}

The following related works provide theoretical and practical solutions to the discounted return case in the presence of various types of delay with the exception of reward collection delays~(\Cref{subsec:reward_delay_related}) which mainly focuses on average rewards. 
We start by describing the three main approaches from the literature to tackle constant delays in state observation or action execution. 
These are the augmented state approach (\Cref{subsec:augmented_related}), the memoryless approach (\Cref{subsec:memoryless}) and the model-based approach (\Cref{subsec:model_based_related_works}).
Then, we extend our overview to include a wider spectrum of delays, still concerning action execution and state observation; research on stochastic delays is presented in \Cref{subsec:stochastic_related} and on non-integer delays in \Cref{subsec:non_int_related}.
Finally, we discuss the reward collection delays in \Cref{subsec:reward_delay_related}.



\subsection{Augmented State Approaches for Constant Delays}
\label{subsec:augmented_related}

\subsubsection{Definition of an Augmented State}
\label{subsubsec:theory_augmented_space}
This section introduced an important concept for \emph{constant} state observation and action execution delays.  
Respectively, \cite{bertsekas1987dynamic} and \cite{altman1992closed} showed that in the case of action execution and state observation delays, the \gls{dmdp} can be reduced to an \gls{mdp} by augmenting the state space as follows. Consider an action execution or state observation delay $\delay$; let $(a_1,\dots,a_{\delay})$
be the last $\delay$ actions that the agent has already selected but whose outcomes it has not yet observed, because of the said delay; let $s$ be the last state the agent has observed. 
Then \emph{augmented state} $x=(s,a_1,\dots,a_\delay)$ contains all the information that the agent can possibly gather at that time. 
Any further information would be redundant under the Markovian property of the underlying \gls{mdp}.
The result of \cite{bertsekas1987dynamic,altman1992closed} is that this new process, obtained by augmenting the state, is indeed an \gls{mdp}.
Its new state now lives in $\augmentedstatespace=\statespace\times\actionspace^\delay$. 
Note the exponential dependency on the delay. 
We note $\delayedpolicyspace$ the set of policies which depends on the augmented state or on the history of augmented states. 
Said alternatively, $\delayedpolicyspace$ contains any policy from $\policyspace$ that does not depend on the information that is more up-to-date than the information contained in the augmented state.
The augmented state and how the policy is based on it is represented in \Cref{fig:state_obs_delayed_mdp} and \Cref{fig:action_exec_delayed_mdp}.
Let us now characterise this process. For some augmented states $x,x'\in\augmentedstatespace$, such that $x=(s,a_1,\dots,a_{\delay})$ and $x'=(s',a_1',\dots,a_{\delay}')$ and for some action $a\in\actionspace$, the  
new transition function reads\footnote{The notation "~~$\widetilde{}$~~" will be used to refer to delayed quantities.}
\begin{align*}
    \augmentedtransitionfunction(x'\vert x,a) = p(s'\vert s,a_1)\delta_a(a_{\delay}')\prod_{i=1}^{\delay-1}\delta_{a_{i+1}}(a_{i}'),
\end{align*}
where the product of Dirac functions ensures that the sequence of actions is correctly passed from an augmented state to the next. 
The reward function reads,
\begin{align}
\label{eq:delayed_reward_deterministic}
    \augmentedexpectedrewardfunction(x,a) = r(s,a_1).
\end{align}
As shown by \cite[Lemma~1]{katsikopoulos2003markov}, an equivalent formulation could be
\begin{align}
\label{eq:delayed_reward_expected}
    \augmentedexpectedrewardfunction(x,a) = \expectedvalue_{z\sim \augmentedbelief(\cdot\vert x)} [r(z,a)],
\end{align}
where $\augmentedbelief$ is the probability of the current unknown state $z$ knowing the augmented state $x$. We refer to this probability as \emph{belief}, for its similarity to the concept of belief in the \gls{pomdp} literature. 
Intuitively, since the current unobserved state distribution is entirely known given the sequence of actions and the last observed state, the first process will eventually collect the same reward in expectation as the second, after $\delay$ steps.
Therefore, the optimal policies in both processes correspond. 
However, note that due to the delay initialisation shift discussed in \Cref{subsec:delay_init}, the two processes can have different returns. 

For completeness, we provide the expression of the belief of the unobserved current state $s_{\delay+1}$ knowing the augmented state $x=(s_1,a_1,\dots.a_\delay)$,
\begin{align}
\label{eq:belief_def}
    \augmentedbelief(s_{\delay+1}\vert x)=\int_{\statespace^{\delay-1}}p(s_{\delay+1}\vert s_\delay,a_\delay)\prod_{i=2}^\delay p(s_i|s_{i-1},a_{i-1})\de s_{2:\delay}.
\end{align}

\subsubsection{Equivalence of Action Execution and State Observation Delays}
\label{subsubsec:equiv_action_state_delays}
Another important result of constant action execution and state observation delays is due to \cite{katsikopoulos2003markov}. 
It states that these two types of delay are two faces of the same coin. 

\begin{prop}[Equivalence of state observation and action execution delays, Result~1 of \cite{katsikopoulos2003markov}]
    A \gls{dmdp} with constant action execution delay $\delay^a$ and constant state observation delay $\delay^s$ can be reduced to an \gls{mdp} by augmenting the state with the last $\delay^s+\delay^a$ actions.
\end{prop}

\subsubsection{Practical Applications}
A relatively straightforward line of research has thus been to consider the direct application of \gls{rl} algorithms to the augmented state \gls{mdp}.
This approach dates back to \cite{brooks1972markov}.
The theory also supports this approach as a way to cast the problem back to an undelayed \gls{mdp}, where the agent can achieve an optimal return for the original \gls{dmdp}. 
However, there is a hindrance; the growth of the augmented state space is exponential in delay since $\lvert\augmentedstatespace\rvert= \lvert\statespace\rvert\lvert\actionspace\rvert^\delay$ \cite{walsh2009learning}. 
This can drastically affect learning speed. 
Nonetheless, more recent works propose revisiting the approach to accelerate learning. 
It is possible, for instance, to modify the actions inside the action buffer of the augmented state in order to simulate the effect of applying a different policy without requiring additional sampling in the environment. 
This ingenious technique, proposed by \cite{bouteiller2020reinforcement}, provides much better sample efficiency. 
Their algorithm, Delay-Correcting Actor-Critic (DCAC) builds on \gls{sac} \cite{haarnoja2018soft}  to which it adds the above idea. 
The experimental results demonstrate very clearly the efficiency of the approach.
The \gls{sac} algorithm seems particularly adapted to the problem of delays, and its core concepts have also been included in other approaches such as RTAC \cite{ramstedt2019real}, an actor-critic algorithm that uses the augmented state as input. 

The augmented state approach is also used in real-time \gls{rl}.
In \cite{xiao2020thinking}, the authors derive theoretical results for continuous-time \gls{rl} where the time required for action selection is taken into account. 
In this setting, too, augmenting the state brings back the properties of a continuous-time undelayed \glspl{mdp}.
Interestingly, the empirical study suggests that extra information--not required in the augmented state and thus in theory--provides higher returns.
Said alternatively, additional information, although theoretically redundant, might make the problem simpler from a learning point of view. 



\subsection{Memoryless Approaches for Constant Delays}
\label{subsec:memoryless}
Inspired by the \gls{pomdp} literature, a second line of search is interested in applying \gls{rl} algorithm to the memoryless process. We refer to this approach as the \emph{memoryless} one.
This process considers only the last observed state and discards the possible information contained in the sequence of actions contained in the augmented state. 
However, note that the approach can still take into account the delay in some way, as dSARSA \cite{schuitema2010control}. 
Interestingly, the authors have been inspired to apply this variation of SARSA \cite{rummery1994line} to a memoryless state by the good results that SARSA has obtained on \glspl{pomdp} problems \cite{loch1998using}. 
The dSARSA algorithm takes into account the delay during the TD update. Let $t$ be the current time.
While the traditional version of SARSA would assign the reward collected at time $t$ to the last selected action $a_t$ and the last observed state $s_{t-\delay}$, dSARSA instead assigns this reward to the same last observed state, but associated with the oldest action in the augmented state $s_{t-\delay}$. 
In fact, this is the action applied at $s_{t-\delay}$. 
This modification has the effect of recovering the definition of the reward of \Cref{eq:delayed_reward_expected}.
The experiments provided by \cite{schuitema2010control} show that 
dSARSA performs well in practice.

In the real-time \gls{rl} literature, \cite{hester2013texplore} consider Monte Carlo Tree Search planning to simulate trajectories given the last observed state. 
Note that the algorithm is specifically designed to reduce the delay in action selection by explicitly separating action selection from the model learning phase.
In fact, the real-time literature \gls{rl} has an important difference from delayed \gls{rl} in that modifying the online algorithm used to learn the policy can reduce the delay resulting from it \cite{hester2012rtmba,caarls2015parallel,ramstedt2019real}. 
Instead, in this dissertation, we consider that the delay is a feature of the environment and that the agent has no impact on it. 

\subsection{Model-Based Approaches for Constant Delays}
\label{subsec:model_based_related_works}
The last approach encountered in the literature is what we will refer to as \emph{model-based}.
The idea is to address the computational cost induced by the exponential dependence of the state on delay $\delay$. To alleviate this cost, these approaches compute, from the augmented state, some substitute for the current unobserved state and use it as input to the policy. 
This way of anticipating the future, or rather the unobserved present, is relatively natural. 
As noted by \cite{firoiu2018human} who cite experimental psychology results, the brain accounts for the potential delay between observation and actuation by extrapolating the future position of moving objects \cite{nijhawan1994motion}.
The aforementioned state substitute typically has a much smaller dimension than the augmented state.
It can be any statistics of the current state, such as its expected value or mode. 
The name model-based stems from the fact that these approaches learn a model of the environment's dynamics in order to recursively simulate the effect of the actions contained in the buffer and obtain some prediction over the current state.

Related works have evaluated learning the most probable current state \cite{walsh2009learning} to use it as input to the policy. In the case of deterministic \gls{mdp}, assuming that the model is perfect, the agent could compute the exact current state and apply the optimal undelayed policy. 
Therefore, the delayed policy thus obtained would match the performance of the undelayed one. 
It is a well-known fact in the delay literature that deterministic \glspl{mdp} induce \glspl{dmdp} that have the same optimal return. 
The stochasticity of the underlying \gls{mdp} is responsible for weaker optimal performance in the delayed case. 
The conjoined effect of stochasticity and delay is well illustrated in \cite[Remark 3.1]{derman2021acting}.
In \cite{walsh2009learning}, the authors extend their algorithm to \emph{mildly stochastic} \glspl{mdp}--\glspl{mdp} where there exists $0\le\epsilon<1$ such that $\forall (s,a)\in\statespace\times\actionspace, \exists s', p(s'\vert s,a)\geq 1-\epsilon$. 
Under this assumption, the expected discounted value function of the delayed policy, $\widetilde{V}^{\pi}$, has the following guarantee \wrt the undelayed policy, $V^{\pi}$, it is based upon: 
\begin{align*}
    \lVert \widetilde{V}^{\pi}-V^{\pi}\rVert_{\infty}\leq \frac{\gamma\epsilon \maximumreward}{(1-\gamma)^2},
\end{align*}
where $\maximumreward$ is the maximum absolute reward.
We note that the assumption of mildly stochastic \gls{mdp} is quite strong and that the bound grows quadratically in the effective horizon $\textstyle\frac{1}{1-\gamma}$.

More recent approaches try to learn the current state using \glspl{nn}, with linear \cite{derman2021acting} or recurrent layers \cite{firoiu2018human}. 
This latter work trains the model of the environment to predict the current state component-wise, using $L2$ loss for continuous elements and the cross-entropy for discrete ones.
The model, therefore, learns to predict the expected value of continuous elements. 
The prediction is then fed to the undelayed \gls{rl} algorithm IMPALA \cite{espeholt2018impala}.
In \cite{derman2021acting}, the predicted state is used as input for Q-learning \cite{watkins1989learning}.  

Using the Q-function as well, the approach of \cite{agarwal2021blind} differs from the previous one in an interesting manner. 
The authors propose to select an action that maximises the expected value of some undelayed Q-function over the distribution of the current unobserved state.
This distribution is given by the belief, which is obtained from the augmented state as in \Cref{eq:belief_def}.
The algorithm, called Expectation Maximization Q-Learning (EMQL), obviously suffers from the fact that the Q-function is computed for an undelayed policy and it is not said that the obtained delayed policy may gather the same reward. 
However, the authors provide theoretical results showing that such a delayed policy $\delayedpolicy$ has the following lower bound guarantee on its value function $\delayedvaluefunction$ taken in some augmented state $x\in\augmentedstatespace$,
\begin{align*}
    \delayedvaluefunction^\delayedpolicy(x) \ge \expectedvalue_{s\sim \augmentedbelief(\cdot\vert x)}\left[V^{*}(s)\right] - \frac{\maximumreward}{(1-\gamma)^2}\left(1-\frac{1}{\cardinal{\actionspace}}\right),
\end{align*}
where $V^{*}$ is the value function of the optimal undelayed policy.

Going further, \cite{firoiu2018human} note that the learnt model could be better used, for instance, to perform planning beyond the current unobserved state.
This is precisely what \cite{chen2021delay} have implemented. 
In their approach, Delay-Aware Trajectory Sampling (DATS),
the model of the \gls{mdp} is learnt as an ensemble
of Gaussian distributions represented by probabilistic neural networks. 
The model is then queried to sample the current unobserved state but also beyond, in order to plan for an action sequence several steps after this current state.  
The approach is based on the probabilistic ensemble
with trajectory sampling (PETS) \cite{chua2018deep}, a model-based approach that has achieved great performances compared to model-free ones such as \gls{ppo}. 
The advantage of PETS in the delayed setting is even clearer.


\subsection{Stochastic Delays}
\label{subsec:stochastic_related}
While most of the aforementioned works assume a constant delay, some also tackle stochastic delays.
\cite{agarwal2021blind} consider delays sampled following a geometric distribution. 
This implies that newer observations may anticipate older ones. 
Consequently, older observations might be collected when fresher information is already available. 
This may be problematic, as the agent loses count of the most up-to-date observation and uses older information instead. 
This phenomenon would lead to the problem of anonymous delay.
To avoid it, \cite{agarwal2021blind} provides the agent with the time stamps of the observations in order to recover the non-anonymity.
Stochastic delays are also studied in \cite{bouteiller2020reinforcement}, where state and action delays follow a bounded discrete distribution. 
As a solution, the authors propose to augment the state to cast the problem back to an \gls{mdp}, similarly to the constant delay case.
However, not only the state is augmented with the last $\delay^a+\delay^s$ actions--to account for the total equivalent delay--but it is also augmented by the current value of both the state delay and the action delay. 
Although this information is not useful in the case of constant delay, here it can be used by the agent to learn about the delay process.
Their DCAC algorithm obtains great empirical results when tested against realistic stochastic delays due to WiFi communication.


\subsection{Non-Integer Delays}
\label{subsec:non_int_related}
To the best of our knowledge, only one work has considered the non-integer delay problem in the \gls{rl} literature. 
In \cite{schuitema2010control}, a paper that has already been discussed in memoryless approaches, the authors show how to adapt their algorithm to this problem.
An important assumption for their results is that the transition of the continuous-time process can be linearised at any state $s$ and for any action $a$, that is, there exist some matrices $B$ and $C$ such that,
\begin{align*}
    \dot{s}(t)=Bs(t) + Ca(t).
\end{align*}
For a delay $\delay\in(0,1)$, at time $t$, the action $a_{t-1}$ is still being applied, and the action that the agent chooses at $t$ will start to be applied at $t+\delay$. 
Then, under the linearisation assumption, the effective action between $t$ and $t+1$ will be $\hat{a}_t = \delay a_t + (1-\delay) a_{t+1}$. 
Therefore, \cite{schuitema2010control} propose to update the Q-function of the pair $(s_t,\hat{a}_t)$ in this case. 
Experimental results show that this modification helps dSARSA learn to reach the goal state faster than SARSA, even if the latter is provided with the augmented state. 


\subsection{Reward Collection Delays}
\label{subsec:reward_delay_related}
As we have anticipated in \Cref{subsec:delay_in_this_thesis}, the problem of delayed reward collection is beyond the scope of this thesis. However, we provide some references for the interested reader. 
A delayed reward has an impact on performance mainly if one is interested in the performance during the learning phase. 
Instead, if one is only interested in the final performance of the agent, having a delayed reward is not highly important as long as it is non-anonymous. 
The reward will eventually be credited to the correct transition. 
All practical offline \gls{rl} algorithms will not be affected by this problem. 

Therefore, it is natural to observe that the delay in reward collection has been mainly discussed in the \emph{online learning} literature.
This field considers an agent who has to choose repeatedly between $K$ arms. 
Each arm provides a reward to the agent which can be either generated in a stochastic or adversarial manner. 
The environment has no state and the agent can always choose from the $K$ arms at any time. 
It is possible to represent this framework as an \gls{mdp} with a unique state, \ie $\cardinal{\statespace}=1$. 
When the agent can only observe the reward of the arm it has chosen and not the reward of the other $K-1$ arms, the setting is called a \emph{multi-armed bandit} \cite{lattimore2020bandit}. 
When the reward is stochastic, the bandit is called \emph{stochastic bandit} while if it is adversarial, it is referred to as \emph{adversarial bandit}.
In this field, an important concept is \emph{regret}, that is, the expected difference between the reward collected by the learning agent and the best policy in hindsight. 

In the bandit literature, \cite{joulani2013online} have shown that, provided that its expected value is finite, the delay has an additive effect on regret in the stochastic bandits setting and a multiplicative effect in the adversarial bandits one.
Their solution is based on the famous concept of Upper Confidence Bounds (UCB) \cite{auer2002finite} and extends it simply by omitting rewards that have not yet been collected.
Many interesting directions have been taken from there. 
For example, \cite{lancewicki2021stochastic} study the effect of a delay that depends on the current reward in stochastic bandits. 
The authors study the case of possibly unbounded support and expectation distributions and derive an algorithm with an additive penalty depending on the quantile of the delay in the regret. 
This bound is also close to the lower bound proved by the authors. 
In their solution, \cite{lancewicki2021stochastic} use \emph{successive elimination}, an algorithm that eliminates arms once the confidence that they are sub-optimal is high enough. Interestingly, they show that this algorithm has greater guarantees than UCB-type algorithms when applied to delays. 
In \cite{romano2022multi}, a new delay setting is proposed.
The reward from pulling an arm can be spread--or delayed--through multiple future steps. 
This property is named \emph{temporally-partitioned rewards} by the authors.
Therefore, the agent observes several partial rewards--or per-round rewards--at each step, coming from multiple past pulls. 
Note that \cite{romano2022multi} consider the case in which the agent is aware of which arm generated each of these partial rewards. 
Said alternatively, the setting is that of non-anonymous delays.
To deal with temporally partitioned rewards, the authors assume a maximum delay and consider that the reward spreads smoothly across the steps until the maximum delay. 
The solution they propose builds an upper confidence bound on the estimated arm's expected reward, where the estimator assumes a 0 per-round reward for unobserved outcomes. 


The branch of \gls{rl} that is interested in the regret of the algorithm during the learning phase, theoretical \gls{rl}, has been interested in the problem of delayed reward collection more recently. 
\cite{howson2021delayed} show results for a wide range of classic algorithms such as UCRL \cite{auer2008near} in stochastic \glspl{mdp}. 
Interestingly, as in \cite{joulani2013online}, they show an additive regret proportional to the delay in most cases. 
In \cite{campbell2016multiple}, the authors use Q-learning to address the problem of a Poisson-distributed reward. 
The idea is to learn multiple Q-function, one for each potential value of the mean delay of the distribution, $\lambda$. 
In this way, each time a new reward is collected, all Q values can be updated as if the reward had been collected from a distribution with the corresponding mean. 
These Q-functions can be seen as a single Q-function where the input is augmented with the value of $\lambda$. 
In this way, the agent can select its next action taking the maximum over all $\lambda$ before selecting the action that maximises this particular Q value. 
The authors prove the convergence of this algorithm. 


\subsection{Anonymous Delays}
\label{subsec:anonymous_delay_related}

First, we discuss anonymous reward collection delays.
In the bandit literature, \cite{pike2018bandits} study the challenging setting of anonymous delays in the stochastic bandit setting. 
If the expectation of the delay is known, the authors propose an algorithm that has an additive term in the delay that also depends logarithmically on $H$, the horizon. 
Interestingly, if the delay is bounded, \cite{pike2018bandits} prove a bound that matches that of \cite{joulani2013online}. 
Their solution repeatedly selects the same arm over a period of time, so that the probability that new rewards come from this arm increases. 
By carefully selecting the length of such a period, \cite{pike2018bandits} are able to derive confidence intervals on the expected reward of an arm. 
This enables the removal of suboptimal arms.

As we have seen in \Cref{subsec:reward_delay}, delayed reward in the literature may refer to credit assignment problems, where all rewards are assigned to a single transition to simplify the design of the environment. 
As an illustration, we report the example provided in \cite{gong2019decentralized} of a traffic congestion reduction problem. 
In this problem, defining a per-step reward for the impact on traffic congestion of a single traffic light is complex. 
However, it is easier to use the average routing time as a reward, but it is only accessible in the final step.
These credit assignment problems can be seen as instances of anonymous delays since the single-step rewards are not clear, and only an aggregated reward is provided in the end.
Many works in the literature deal with this problem, including, for instance, \cite{arjona2019rudder} or \cite{han2022off}. 
Yet, to the best of our knowledge, no work in the \gls{rl} literature considers the problem of well-defined per-step rewards with anonymous delays.

Finally, with respect to the state observation or action execution delays in \gls{rl}, we are not aware of any work.
This could indeed constitute an interesting research direction.

\subsection{Overview of Delays in Reinforcement Learning}
We give an overview of the literature on state observation and action execution delays in \Cref{tab:delay_bestiary}.
In \Cref{tab:delay_related} we provide a similar table but focus on the approaches that have been considered and the delay to which they have been applied. 
This allows one to see where results are missing.

\begin{figure}[t]
% Inspired from https://tex.stackexchange.com/questions/78938/tables-in-latex-with-diagonal-first-row
% and https://superuser.com/questions/1148690/how-can-i-create-a-sophisticated-table-like-the-one-attached
    \centering
    \begin{tikzpicture}[remember picture]
    \node[inner xsep=-\pgflinewidth,inner ysep=-\pgflinewidth] at (0,0) (mytable){%
    \begin{tabular}{L{c00}|T{c0}T{c1}T{c2}T{c3}T{c4}T{c5}T{c6}T{c7}T{c8}||T{c9}T{c10}}
        \cline{1-11}\\
        \cite{brooks1972markov} & \checkmark &-& \checkmark &-& $1$ & \checkmark &-&-&\checkmark&A\\
        \cite{walsh2009learning} & \checkmark &-& \checkmark &-& $1-20$ & \checkmark &-&-&\checkmark&Mo\\
        \cite{schuitema2010control} & \checkmark &-& \checkmark &-& $0-2$ & \checkmark &\checkmark&-&\checkmark&Me\\
        \cite{hester2013texplore} & \checkmark &-& \checkmark &-& $2$ & \checkmark &-&-&\checkmark&Me\\
        \cite{firoiu2018human}& \checkmark &-& \checkmark &-& $1-7$ & \checkmark &-&-&\checkmark&Mo\\
        \cite{ramstedt2019real} & \checkmark &-& \checkmark &-& $1$ & \checkmark &-&-&\checkmark&A\\
        \cite{bouteiller2020reinforcement} & \checkmark &-& \checkmark&\checkmark& $1-5$ & \checkmark &-&-&\checkmark&A\\
        \cite{agarwal2021blind} & \checkmark &-& \checkmark &\checkmark& $2-4$ & \checkmark &-&-&\checkmark&Mo\\
        \cite{derman2021acting} & \checkmark &-& \checkmark &-& $1-25$ & \checkmark &-&-&\checkmark&Mo\\
        %\cite{liotet2021learning} & \checkmark &-& \checkmark &-& $3-15$ & \checkmark &-&-&\checkmark&Mo\\
        \cite{chen2021delay} & \checkmark &-& \checkmark &-& $1-16$ & \checkmark &-&-&\checkmark&Mo\\
        %\cite{liotet2022delayed} & \checkmark &-& \checkmark &-& $0-20$ & \checkmark & \checkmark &-&\checkmark&\\
    \end{tabular}
    };
    
    % Diagonal separations
    \coordinate (A) at ($(c8)+(1mm,0)$);
    \draw (mytable.north-|c1) --++ (60:3cm);
    \draw (mytable.north-|c3) --++ (60:3cm);
    \draw (mytable.north-|c4) --++ (60:3cm);
    \draw (mytable.north-|c6) --++ (60:3cm);
    \draw (mytable.north-|c8) --++ (60:3cm);
    \draw (mytable.north-|A) --++ (60:3cm);
    
    % Column names
    \node[rotate=60,anchor=west] at ($(mytable.north-|c00)!0.5!(mytable.north-|c0)$) {State/Action};
    \node[rotate=60,anchor=west] at ($(mytable.north-|c0)!0.5!(mytable.north-|c1)$) {Reward};
    \node[rotate=60,anchor=west] at ($(mytable.north-|c1)!0.5!(mytable.north-|c2)$) {Constant};
    \node[rotate=60,anchor=west] at ($(mytable.north-|c2)!0.5!(mytable.north-|c3)$) {Stochastic};
    
    \node[rotate=60,anchor=west] at ($(mytable.north-|c3)!0.5!(mytable.north-|c4)$) {Amount of Delay};
    \node[rotate=60,anchor=west] at ($(mytable.north-|c4)!0.5!(mytable.north-|c5)$) {Integer};
    \node[rotate=60,anchor=west] at ($(mytable.north-|c5)!0.5!(mytable.north-|c6)$) {Non-integer};
    \node[rotate=60,anchor=west] at ($(mytable.north-|c6)!0.5!(mytable.north-|c7)$) {Anonymous};
    \node[rotate=60,anchor=west] at ($(mytable.north-|c7)!0.5!(mytable.north-|c8)$) {Non-Anonymous};
    \node[rotate=60,anchor=west] at ($(mytable.north-|c8)!0.5!(mytable.north-|c9)$) {Approach};
    
    \end{tikzpicture}
    \captionof{table}{Related works and the type of delay that they consider. 
    The ordering is with respect to the publication date.
    The amount of delay corresponds to the range of delay that has been taken into account in the experiments. 
    For stochastic delays, we only give the mean delay as the delays can potentially be infinite.
    The different approaches are labelled $\textbf{A}$ for the augmented approach, $\textbf{Me}$ for the memoryless one and  $\textbf{Mo}$ for the model-based one.}
    \label{tab:delay_bestiary}
\end{figure}


\begin{table}[t]
\begin{tabular}{lll||>{\raggedright\arraybackslash}p{2.1cm}|>{\raggedright\arraybackslash}p{2.1cm}|>{\raggedright\arraybackslash}p{2.1cm}}

 \multicolumn{3}{c||}{Type of Delay}&  Augmented & Memoryless & Model-based \\
 \hline\hline
Constant & Integer & Non-anonymous&  {\footnotesize\cite{brooks1972markov}\cite{ramstedt2019real}\cite{bouteiller2020reinforcement}} &  {\footnotesize\cite{schuitema2010control}\cite{hester2013texplore}} & {\footnotesize\cite{walsh2009learning}\cite{firoiu2018human}\cite{agarwal2021blind}\cite{derman2021acting} \cite{chen2021delay}}\\
 \hline
Stochastic & Integer & Non-anonymous& {\footnotesize\cite{bouteiller2020reinforcement}} &  & {\footnotesize\cite{agarwal2021blind}}  \\
&&&&&\\
 &&&&&\\
\hline
Constant & Non-integer & Non-anonymous&  & {\footnotesize\cite{schuitema2010control}} &  \\
&&&&&\\
&&&&&\\
 \hline
Constant & Non-integer & Non-anonymous&  &  &  \\
&&&&&\\
&&&&&\\
\end{tabular}
\caption{Overview of how the literature has addressed different types of state observation and action execution delays.}
\label{tab:delay_related}
\end{table}


\subsection{Delays in Control Theory}
\label{subsec:control_theory}
We finish this section with what we believe to be an important extension of the review of the literature. 
Control theory shares much in common with \gls{rl}, both frameworks consider a decision-making process, and control theory also considers the problem of delays, as we have anticipated in \Cref{sec:why_delay_ns}.
Therefore, it could be useful to leverage the results and ideas of this framework to apply them in \gls{rl}.
This is even more true since control theory has a long history and the delay that it has considered is typically much more diverse. 
For example, the case of a delay that, although constant, is different for each dimension of the state of the environment \cite{alevisakis1973extension, altman1999congestion}. 
Interestingly, in this case, as in \gls{rl}, the state observation and action execution delays have been shown to be equivalent.
Another type of delay is distributed delay, which appears when the system depends on a continuous interval of past states \cite{gu2003survey}. 
It is more common to work with continuous-time systems in control theory than in \gls{rl}.
For classical constant delays, as for \gls{rl}, the state can be augmented to cast the delayed system back to an undelayed one \cite{kwon2003simple}.
Another approach for constant delays is the Smith Predictor model \cite{smith1957closer}.
It proposes to learn a model of the undelayed process and leverage this model to cast the problem back to the undelayed problem. 
A link can be drawn between this approach and the aforementioned model-based approaches to constant delays in \gls{rl}.

Another layer of complexity can be added to the problem by considering the delay in a multi-agent framework, that is, when several agents can act simultaneously--or not--in the environment.
The problem then becomes \emph{decentralized}, and since each agent is potentially independent, each can have a different delay.
A model in the control literature for such a system is \emph{delayed sharing information pattern}.
As defined by \cite{witsenhausen1971separation}, the problem considers a set of $K$ agents, each with a state $s_t^k$ that composes the state $\boldsymbol s_t$ of the environment at time $t$.
Each agent selects an action $a_t^k$ that composes the global action $\boldsymbol a_t$ applied to the environment.
As in \gls{pomdp} (\Cref{subsec:pomdp}), the agents do not directly observe the state of the environment but only a partial observation $\boldsymbol o_t = (o_t^k)_{1\le k \le K}$ of it. 
Each of these agents is $n$-delayed in the sense that it receives information from the other agents with a delay $n$ while it has undelayed access to its own state.
Therefore, at time $t$, agent $k$ observes $o_{h}^k$ for $h\le t$ but observes only $o_{h}^j$ for $j\ne k$ and $h\le t-n$.
Two key concepts of this framework are the information shared by all agents at time $t$, which reads,
\begin{align*}
    A_t = \left\{ (\boldsymbol o_h, \boldsymbol a_h)\vert h\le t-n \right\},
\end{align*}
and the additional information available to an agent $k$ at time $t$ that reads,
\begin{align*}
    B_t^k = \left\{ ( o_h^k,  a_h^k)\vert t-n<h\le t-1 \right\} \cup \{ o_t^k\}.
\end{align*}
Note that at time $t$ the agent does not have access to the action $ a_t^k$, hence the last term of the previous equation.
\cite{witsenhausen1971separation} conjectured that, instead of considering all policies based on the information in $(A_t,B_t^1,\dots,B_t^k)$, there existed an optimal policy where the belief of the current state $\boldsymbol s_t$ given $A_t$ would be substituted for the term $A_t$ in the previous information structure. 
This is a similar idea to what is proposed by model-based delayed \gls{rl} algorithms. 
In control theory, for delays greater than 1, the conjecture has been shown false by \cite{varaiya1978delayed}.
In \gls{rl}, a similar result holds, as we will see in \Cref{sec:th_analysis_belief}.
Another similarity with delayed \gls{rl} is the following.
When $n=0$--i.e. the state itself is observable--the problem can be formulated as an \gls{mdp} and, as for the augmented state approach, it has been shown that the problem of 1-delayed sharing information pattern could be cast back to an \gls{mdp} by augmenting the state with the delayed action \cite{hsu1982decentralized}.
However, contrary to delayed \gls{rl}, some results show that not all history is necessary to design optimal policies, and reduced information structures consisting only of some statistics that do not grow with time can be used instead \cite{altman2009stochastic,nayyar2010optimal}.

Despite the similarities between control theory and \gls{rl}, the delay problems considered in the control theory literature are usually more complex and therefore more involved from a theoretical perspective. 
Instead, the \gls{rl} framework and its \gls{mdp} process offer a simpler framework where structural results can be more easily derived.  